% Section 1 - Templates
% Roberto Masocco <roberto.masocco@uniroma2.it>
% April 22, 2026

% ### Templates ###
\section{Templates}
\graphicspath{{figs/section1/}}

% --- Templates ---
\begin{frame}{Templates}{Generic programming}
	\justifying
	Consider writing a \texttt{max()} \textbg{function} for \textbg{integers}, then again for \textbg{doubles}, then again for \textbg{any type} that supports comparison.\\
	In \textbg{C} the only escape is a \textbg{macro} — unsafe, untyped, and undebuggable.\\
	\bigskip
	C++ solves this with \textbg{templates}: a mechanism for writing code that is \textbg{parameterized on a type} (or on a value).\\
	\textbg{The compiler generates the actual code} for each combination of template arguments used in a program — at \textbg{compile time}, with full type checking.\\
	\bigskip
	\begin{block}{Definition of template}
		A \textbf{template} is a blueprint for a \textbf{family} of functions or classes, parameterized by one or more \textbf{type parameters} (or non-type parameters), from which \textbf{the compiler generates concrete instantiations} on demand.
	\end{block}
\end{frame}
\begin{frame}[fragile]{Templates}{Function templates}
	\justifying
	A \textbg{function template} is declared by prefixing the function with a \textbg{template parameter list}.\\
	The type parameter is conventionally named \texttt{T} (though any name works).
	\begin{columns}
		\column{.9\textwidth}
		\begin{lstlisting}[style=codebeamer, language=C++, caption=A function template and some instantiations.]
template<typename T>
T max(T a, T b)
{
  return (a > b) ? a : b;
}

int    m1 = max<int>(3, 7);          // T = int
double m2 = max<double>(3.14, 2.71); // T = double
int    m3 = max(10, 20);             // T deduced by compiler
    \end{lstlisting}
	\end{columns}
	When the arguments allow it, \textbg{the compiler deduces} the template argument automatically.
\end{frame}
\begin{frame}[fragile]{Templates}{Class templates}
	\justifying
	A \textbg{class template} defines a family of classes.\\
	All member functions of a class template are \textbg{implicitly templates} themselves.
	\begin{columns}
		\column{.9\textwidth}
		\begin{lstlisting}[style=codebeamer, language=C++, caption=A minimal class template.]
template<typename T>
class Box
{
public:
  Box(T value) : value_(value) {}
  T get() const { return value_; }
private:
  T value_;
};
Box<int>         int_box(42);
Box<std::string> str_box("hello");
    \end{lstlisting}
	\end{columns}
\end{frame}
\begin{frame}[fragile]{Templates}{Instantiation and specialization}
	\justifying
	Each distinct set of template arguments produces a \textbg{separate instantiation} — a completely independent piece of compiled code.\\
	A \textbg{template specialization} overrides the generic definition for that type only.
	\begin{columns}
		\column{.9\textwidth}
		\begin{lstlisting}[style=codebeamer, language=C++, caption=Full specialization of \texttt{Box} for \texttt{bool}.]
template<>
class Box<bool>
{
public:
  Box(bool value) : value_(value) {}
  bool get() const { return value_; }
  void flip() { value_ = !value_; }
private:
  bool value_;
};
    \end{lstlisting}
	\end{columns}
\end{frame}
\begin{frame}[fragile]{Templates}{Taming instantiation names}
	\justifying
	Template instantiation names grow quickly and become unwieldy in real code.\\
	\texttt{using} aliases give them shorter, meaningful names.
	\begin{columns}
		\column{.9\textwidth}
		\begin{lstlisting}[style=codebeamer, language=C++, caption=Type aliases for template instantiations.]
// Instead of writing this everywhere:
std::map<std::string, std::vector<int>> table;

// Define an alias:
using StringIntVecMap = std::map<std::string, std::vector<int>>;
StringIntVecMap table;

// Alias templates are also possible:
template<typename T>
using Vec = std::vector<T>;
Vec<double> data;
    \end{lstlisting}
	\end{columns}
	This pattern is pervasive in ROS 2 — you will see \texttt{using SharedPtr = ...} in virtually every class.
\end{frame}
